\chapter{Introducción y Objetivos}

\section{Objetivos}
\begin{itemize}
    \item Trazar la curva de contrastación de un instrumento a los fines de eliminar el error sistemático del mismo.
    \item Verificar el índice de Clase del instrumento contrastado.
    \item Determinar el valor de la resistencia en situaciones específicas (bajos valores y puestas a tierra).
    \item Dar el resultado de las mediciones acompañadas del grado de incertidumbre.
\end{itemize}

\section{Introducción Teórica}
Todo instrumento de medición presenta discrepancias entre su escala y el verdadero valor de la magnitud, debidas a tolerancias de fabricación y alinealidades. En la industria, se elaboran especificaciones estadísticas (Clase o Exactitud) para acotar este margen. Mediante la \textbf{contrastación} contra un instrumento patrón de mayor jerarquía, es posible aislar y eliminar el \textit{error sistemático}, construyendo una \textit{curva de corrección}. En este trabajo, el patrón deberá ser al menos cinco veces mejor que el instrumento bajo prueba.

Además, se estudian métodos especiales para medir \textbf{resistencias de pequeño valor}, donde el uso de un óhmetro convencional falla debido a que las \textit{resistencias de contacto} y de las puntas de prueba se suman al mensurando. El \textbf{método de 4 terminales} soluciona esto inyectando una corriente constante e independizando el circuito de medición de tensión. 
